\documentclass[aspectratio=169,11pt]{beamer}
\def\courseweek{9}
% Shared style for the Introduction to Information Geometry lecture slides.
% Compile the lecture .tex files from this directory.
\providecommand{\courseweek}{1}
\usepackage[T1]{fontenc}
\usepackage[utf8]{inputenc}
\usepackage[sc]{mathpazo}
\renewcommand{\sfdefault}{lmss}
\renewcommand{\ttdefault}{lmtt}
\usepackage{amsmath,amssymb,mathtools}
\usepackage{tikz}
\usetikzlibrary{arrows.meta,calc,positioning}
\usepackage{pgfplots}
\pgfplotsset{compat=1.18}
\usepackage{booktabs}

\definecolor{courseink}{HTML}{173749}
\definecolor{coursemuted}{HTML}{526674}
\definecolor{courseteal}{HTML}{0E7287}
\definecolor{coursecopper}{HTML}{B66B43}
\definecolor{courseline}{HTML}{D4E0E3}
\definecolor{coursepale}{HTML}{EFF5F5}

\setbeamercolor{normal text}{fg=courseink,bg=white}
\setbeamercolor{structure}{fg=courseteal}
\setbeamercolor{frametitle}{fg=courseink,bg=white}
\setbeamercolor{title}{fg=courseink}
\setbeamercolor{subtitle}{fg=coursemuted}
\setbeamercolor{block title}{fg=courseink,bg=white}
\setbeamercolor{block body}{fg=courseink,bg=white}
\setbeamercolor{alerted text}{fg=coursecopper}
\setbeamerfont{title}{family=\rmfamily,series=\bfseries,size=\fontsize{32}{37}}
\setbeamerfont{frametitle}{family=\rmfamily,series=\bfseries,size=\fontsize{26}{31}}
\setbeamerfont{normal text}{size=\fontsize{18}{22}}
\setbeamerfont{subtitle}{size=\fontsize{17}{21}}
\setbeamersize{text margin left=0.78cm,text margin right=0.78cm}
\setbeamertemplate{navigation symbols}{}
\setbeamertemplate{itemize items}[circle]
\setbeamertemplate{enumerate items}[default]
\setbeamertemplate{blocks}[default]
\setlength{\parskip}{0.22em}
\setbeamertemplate{footline}{%
  \leavevmode\hbox to \paperwidth{\hspace{.78cm}%
  {\color{coursemuted}\fontsize{8.5}{10}\selectfont Introduction to Information Geometry \hspace{.32cm} / \hspace{.32cm} Week \courseweek}%
  \hfill{\color{coursemuted}\fontsize{8.5}{10}\selectfont \insertframenumber\hspace{.17cm}/\hspace{.17cm}\inserttotalframenumber}\hspace{.78cm}}%
  \vskip.17cm
}
\setbeamertemplate{frametitle}{%
  \vspace{.13cm}
  \insertframetitle\par
  \vspace{.08cm}
}
\setbeamertemplate{title page}{%
  \vspace*{1.00cm}
  {\color{courseteal}\fontsize{10}{13}\selectfont\bfseries
  INTRODUCTION TO INFORMATION GEOMETRY\par}
  \vspace{.40cm}
  {\usebeamerfont{title}\usebeamercolor[fg]{title}\inserttitle\par}
  \vspace{.22cm}
  {\usebeamerfont{subtitle}\usebeamercolor[fg]{subtitle}\insertsubtitle\par}
  \vfill
  {\color{coursemuted}\fontsize{13}{16}\selectfont Christopher Lee\par}
  \vspace{.34cm}
}

\pgfplotsset{
  courseplot/.style={
    width=11.05cm,
    height=5.15cm,
    axis line style={coursemuted!75},
    tick style={coursemuted!65},
    ticklabel style={font=\scriptsize,color=coursemuted},
    label style={font=\normalsize,color=courseink},
    grid=major,
    major grid style={courseline!70,line width=.4pt},
    line width=1.45pt,
    legend style={font=\small,draw=none,fill=none,at={(.98,.03)},anchor=south east},
    clip=true
  }
}

\newcommand{\highlightmath}[1]{{\color{courseteal}#1}}
\newcommand{\thinrule}{\par\vspace{.25cm}{\color{courseline}\rule{\linewidth}{.8pt}}\par}

% A one-dimensional state space.  x=0 and x=1 are placed at 0 and 9 units.
\newcommand{\stateinterval}[1]{%
 \begin{tikzpicture}[x=1cm,y=1cm]
   \draw[courseink,line width=1.25pt] (0,0)--(9,0);
   \foreach \pos/\lab in {0/0,9/1} {
     \fill[courseink] (\pos,0) circle(2.7pt);
     \node[below=9pt,text=coursemuted] at (\pos,0) {\(\lab\)};
   }
   #1
 \end{tikzpicture}%
}

% Paired pages are separate frames in the PDF, with matching titles.
% Supply the whole frame body for each version; the instructor may edit them
% independently without changing overlay logic.
\newcommand{\coursepair}[3]{%
  \begin{frame}{#1}#2\end{frame}%
  \begin{frame}{#1}#3\end{frame}%
}

% Equilateral representation of the categorical simplex, in barycentric
% coordinates with vertices A=(0,0), B=(1,0), C=(1/2,sqrt(3)/2).
\newcommand{\coursesimplex}[1]{%
 \begin{tikzpicture}[scale=4.9]
  \fill[coursepale] (0,0)--(1,0)--(.5,.8660254)--cycle;
  \draw[courseink,line width=.65pt] (0,0)--(1,0)--(.5,.8660254)--cycle;
  \node[below left,font=\small] at (0,0) {A};
  \node[below right,font=\small] at (1,0) {B};
  \node[above,font=\small] at (.5,.8660254) {C};
  #1
 \end{tikzpicture}%
}

\title{Natural Gradients and Replicator Flows}
\subtitle{Week 9 \(\cdot\) Lecture 1}
\date{}
\begin{document}
\begin{frame}[plain]\titlepage\end{frame}

\begin{frame}{Which direction increases a function fastest?}
For a function \(\Phi:\mathbb R^n\to\mathbb R\), the ordinary gradient
is defined by
\[
d\Phi_p(v)=\nabla\Phi(p)\cdot v.
\]
\vspace{.45cm}
This dot product measures length using Euclidean geometry.
\vspace{.48cm}
Probability distributions come with a different way to measure
the length of a small change.
\end{frame}

\begin{frame}{A gradient depends on the metric}
Given a metric \(g\), define the \emph{metric gradient} by
\[
\highlightmath{g_p(\operatorname{grad}_g\Phi,v)=d\Phi_p(v)}
\quad\text{for every tangent vector }v.
\]
\vspace{.55cm}
In local coordinates, write \(g(v,w)=v^TG(p)w\). Then
\[
\highlightmath{\operatorname{grad}_g\Phi=G(p)^{-1}\nabla\Phi}.
\]
\end{frame}

\begin{frame}{Fisher geometry on a Bernoulli family}
For \(P=(x,1-x)\), the Fisher line element is
\[
ds_F^2=\frac{dx^2}{x(1-x)}.
\]
\vspace{.45cm}
Its metric coefficient is \(G(x)=1/[x(1-x)]\).
\vspace{.50cm}
Consequently the Fisher-gradient ascent equation is
\[
\highlightmath{x'=x(1-x)\Phi'(x)}.
\]
\end{frame}

\coursepair{An ascent flow on the Bernoulli interval}{
Take the potential \(\Phi(x)=3x-2x^2\).
\vspace{.48cm}
Find its natural-gradient equation and any interior equilibrium.
\[
x'=x(1-x)\Phi'(x).
\]
}{
Take the potential \(\Phi(x)=3x-2x^2\).
\vspace{.38cm}
Since \(\Phi'(x)=3-4x\),
\[
\highlightmath{x'=x(1-x)(3-4x)}.
\]
\vspace{.23cm}
The interior equilibrium is \(\highlightmath{x_*=3/4}\).
}

\begin{frame}{The phase line of this flow}
\vspace{.35cm}
\begin{center}
\stateinterval{
  \fill[coursecopper] (6.75,0)circle(3pt);
  \node[below=9pt,text=coursecopper] at (6.75,0) {\(3/4\)};
  \draw[courseteal,-{Stealth[length=4mm]},line width=1.9pt](1.7,.18)--(3.8,.18);
  \draw[courseteal,-{Stealth[length=4mm]},line width=1.9pt](8.45,.18)--(7.35,.18);
}
\end{center}
\vspace{.63cm}
The potential rises until \(x=3/4\) and falls after it.
The gradient flow moves toward the interior maximum.
\end{frame}

\begin{frame}{Why the potential increases}
Along a solution \(x(t)\), the chain rule gives
\[
\frac{d\Phi}{dt}=\Phi'(x)x'.
\]
\vspace{.44cm}
Substitute the Fisher-gradient equation:
\[
\highlightmath{\frac{d\Phi}{dt}=x(1-x)[\Phi'(x)]^2\geq0}.
\]
\vspace{.43cm}
The identity expresses ascent without requiring us to solve the ODE.
\end{frame}

\coursepair{How fast is the potential increasing?}{
At \(x=1/2\), for \(\Phi(x)=3x-2x^2\), calculate
\[
\frac{d\Phi}{dt}=x(1-x)(3-4x)^2.
\]
\vspace{.5cm}
Also find the maximal value of \(\Phi\) on \([0,1]\).
}{
At \(x=1/2\), we have \(3-4x=1\), so
\[
\highlightmath{\frac{d\Phi}{dt}=\tfrac12\cdot\tfrac12\cdot1^2=\tfrac14}.
\]
\vspace{.34cm}
The maximum occurs at \(x=3/4\):
\[
\highlightmath{\Phi(3/4)=\tfrac98}.
\]
}

\begin{frame}{Tangent vectors on a simplex}
For \(p=(p_1,\ldots,p_n)\), the coordinates satisfy
\[\sum_i p_i=1.\]
\vspace{.37cm}
Every tangent vector \(v\) therefore satisfies \(\sum_i v_i=0\).
\vspace{.43cm}
The categorical Fisher metric is
\[
\highlightmath{g_p(v,w)=\sum_i\frac{v_iw_i}{p_i}}.
\]
\end{frame}

\begin{frame}{The Fisher gradient on the simplex}
Let \(a_i=\partial\Phi/\partial p_i\), and set
\(\bar a=\sum_jp_j a_j\).
\vspace{.48cm}
The Fisher gradient has coordinates
\[
\highlightmath{(\operatorname{grad}_F\Phi)_i
=p_i(a_i-\bar a)}.
\]
\vspace{.40cm}
Its coordinates sum to zero, so the resulting velocity is tangent
to the probability simplex.
\end{frame}

\coursepair{A gradient at the center of the simplex}{
Let \(\Phi(p)=2p_A+p_B\), and take
\(P=(1/3,1/3,1/3)\).
\vspace{.46cm}
Find \(\bar a\) and the Fisher-gradient velocity
\[
p_i'=p_i(a_i-\bar a),\quad a=(2,1,0).
\]
}{
Let \(\Phi(p)=2p_A+p_B\), and take
\(P=(1/3,1/3,1/3)\).
\vspace{.36cm}
The weighted mean is \(\bar a=(2+1+0)/3=1\). Hence
\[
\highlightmath{p'=(1/3,\ 0,\ -1/3)}.
\]
\vspace{.28cm}
Type A increases, type B is momentarily constant,
and type C decreases.
}

\begin{frame}{Replicator dynamics is a Fisher gradient}
If the payoff or fitness values have the form
\[
f_i(p)=\frac{\partial\Phi}{\partial p_i},
\]
\vspace{.37cm}
then the natural-gradient flow becomes
\[
\highlightmath{p_i'=p_i(f_i-\bar f)}.
\]
\vspace{.46cm}
This is exactly the replicator equation, now understood as motion
up a potential in Fisher geometry.
\end{frame}

\begin{frame}{Variance measures both speed and ascent}
For the flow \(p'=\operatorname{grad}_F\Phi\),
\[
\frac{d\Phi}{dt}=d\Phi_p(p')=g_p(p',p').
\]
\vspace{.38cm}
Using \(p_i'=p_i(f_i-\bar f)\) gives
\[
\highlightmath{\frac{d\Phi}{dt}=\|p'\|_F^2
=\operatorname{Var}_p(f)}.
\]
\vspace{.32cm}
Fitness variance is the squared geometric speed of ascent.
\end{frame}

\coursepair{Check the ascent identity numerically}{
For \(P=(1/3,1/3,1/3)\) and \(\Phi=2p_A+p_B\),
calculate the instantaneous increase in \(\Phi\).
\vspace{.44cm}
Use \(p'=(1/3,0,-1/3)\) and
\(\dot\Phi=2p_A'+p_B'\).
}{
For \(P=(1/3,1/3,1/3)\) and \(\Phi=2p_A+p_B\),
\[
\highlightmath{\dot\Phi=2(1/3)+0=2/3}.
\]
\vspace{.31cm}
The fitness variance gives the same answer:
\[
\tfrac13(2-1)^2+\tfrac13(1-1)^2+\tfrac13(0-1)^2
=\highlightmath{\tfrac23}.
\]
}

\begin{frame}{The boundary and the metric}
The formula \(p_i'=p_i(f_i-\bar f)\) makes each face
\(p_i=0\) invariant.
\vspace{.53cm}
The Fisher metric becomes singular in these coordinates at that face.
\vspace{.53cm}
Interior trajectories remain in the positive simplex at every
finite time when the fitness functions are sufficiently regular.
\end{frame}

\begin{frame}{A question for the next lecture}
We have described motion that increases a chosen potential.
\vspace{.45cm}
Which potential should we choose if we want a population of
probabilities to approach a specified target distribution?
\vspace{.45cm}
Relative entropy supplies two instructive answers.
\end{frame}
\end{document}
